\chapter{검증 결과(5): R2 보호와 고장 강등}

\section{측정 기반 보호 상태기계(E17)}

창 조정 모듈의 H-브리지 고장은 바로 바이패스하면 수백 A가 흘러(7.2절) 강등 방식이 먼저 정해져야 했다. 실제 장치는 고장을 측정으로 판단하고 판단·차단에 지연이 있으므로, 평균모델에 버스 RC 감쇠기, 링크 제동 저항, 바이패스 소프트 투입(선형 영역 구동), 400 V 부하 분리를 넣고 측정 기반 상태기계를 결합했다. 링크전류 센서 감시(P6)를 켠 기존 시험 격자도 다시 판정했다(1,080건, 판정 변화·오경보 0).

H-브리지 고장의 강등 순서는 H-브리지가 꺼졌을 때 생기는 성질을 쓴다. 꺼진 소자의 역방향 도통이 다이오드 정류기를 이루어 직렬전압이 −(v\textsubscript{link} + 2V\textsubscript{D})로 묶이므로, DAB로 링크전압을 조절하면 버스를 붙잡을 수 있다(그림 \ref{f:seq}).

\begin{figure}[!htb]
\centering
\includegraphics[width=15.5cm]{\repo/docs/fig/report/fig9_1_fault_sequence_final.png}
\caption{H-브리지 고장 강등 순서}\label{f:seq}
\srcnote{주: 9.3–9.4절의 고장 모드 변조, 9.5절의 SALS 연동, 9.6절의 Stage C 일시 정지와 400 V 부하 정지를 포함한 최종 순서다.}
\end{figure}

{\footnotesize
\begin{longtable}{L{\dimexpr 0.4035\linewidth-2\tabcolsep\relax}L{\dimexpr 0.2982\linewidth-2\tabcolsep\relax}L{\dimexpr 0.2982\linewidth-2\tabcolsep\relax}}
\caption{보호 상태기계의 정책 비교와 추천 순서(평균모델)}\label{t:e17}\\
\toprule
\tabhead \textbf{시험} & \textbf{R2-W460} & \textbf{R2-W410} \\
\midrule
\endfirsthead
\toprule
\tabhead \textbf{시험} & \textbf{R2-W460} & \textbf{R2-W410} \\
\midrule
\endhead
\bottomrule
\endlastfoot
H-브리지 개방, 즉시 바이패스 & 0/4, 최대 766 A & 0/4, 최대 869 A \\
H-브리지 개방, 클램프 유지 + 제동(추천) & 4/4, 경로 ≤ 90.0 A, 연결 중 버스 ≤ 467.1 V & 4/4, 경로 ≤ 90.3 A, 버스 ≤ 417.1 V \\
DAB 고장, 제동 유지 → 분리 → 램프 → 바이패스(추천) & 8/8, 링크 ≤ 110.1 V, 제동 ≤ 49.6 J & 8/8, 제동 ≤ 63.5 J \\
켬(전류 조향)·끔(직렬전압 조향) & 켬 6/6·끔 4/4 & 켬 6/6·끔 4/4 \\
추천 순서 전체 & 44/44(켬 12, 끔 8, DAB 고장 16, H-브리지 고장 8) &  \\
H-브리지 변조율 고착(±0.95) & 0/16, 필터 128–147 A(비교기 지연 1.2 µs) &  \\
\end{longtable}
}
\srcnote{주: 추천 순서 전체는 두 모듈을 합한 건수다.}

추천 순서(표 \ref{t:e17})의 켬 순서에서는 바이패스 전류를 H-브리지 영벡터로 먼저 옮긴 뒤(전류 조향) 바이패스를 열어, TB08에서 본 열림 과도를 없앴다(경로 70.3·78.9 A → 57.0·63.8 A). H-브리지 변조율 고착은 비교기 지연 동안 필터전류가 128–147 A까지 올라, 필터 인덕터 포화 전류 ≥ 150 A 또는 더 빠른 비교기가 필요하다.

\section{보호 상태기계의 회로 확인(TB09)}

TB08 회로에 소자 역방향 도통(2 V), 링크 제동 저항(3 Ω), 바이패스 소프트 투입 가변 저항, 버스 RC 감쇠기, 분리 부하, E17과 같은 논리의 상태기계와 보호 게이트(하드웨어 비교기 래치, 고장 주입)를 더했다. 사례는 1008·930 V H-브리지 개방, 930 V 켬, 1008 V DAB 고장이다(모두 25 kW).

\textbf{첫 실행에서는 H-브리지 고장 2사례가 불통과였다.} 전압비가 극단(n·v\textsubscript{link} ≪ v\textsubscript{bus})인 DAB는 링크를 0.5 V까지 내리지 못하고 3.5–22 V에서 멈췄고, 남은 직렬전압을 로그 램프 저항으로 닫자 버스 감쇠기 충전 전류로 경로전류가 177 A까지 올랐다. 붕괴 정체 감지와 경로전류 조절 소프트 투입을 상태기계에 넣어 \textbf{최종 실행 4/4가 통과}했다. 상태 전이 시각은 E17과 0–20 µs 안이었고, 다이오드 클램프 식은 평균 −0.011·−0.050 V 오차로 성립했다.

회로에서만 드러난 스트레스가 하나 남았다. \textbf{강등 중 DAB 순환전류가 정격점의 4.2배}였다(링크 쪽 순시 192 A, 1차 41.7 A, 정격점 45.8 A·약 10 A; SPS 식 계산 197 A). 누설 인덕터를 정격점에 맞춰 설계하면 포화할 수 있고, 4배 정격으로 키우면 창 조정 모듈을 4–5 kW로 줄인 이점이 사라진다.

\section{DAB 고장 모드 변조(E20, TB10)}

원인은 SPS의 전압비 불일치다(그림 \ref{f:e20}). 링크전압이 1차 환산 전압(v\textsubscript{bus}/n)보다 크게 낮으면 전달 전류와 상관없이 누설 인덕터에 큰 전압이 반주기 동안 걸려, 링크 쪽 첨두가 v\textsubscript{bus}/(4·n·f\textsubscript{s}·L′)에 가까워진다(1008 V 팩에서 약 197 A, L′ = L\textsubscript{dab}/n\textsuperscript{2}). 강등 순서를 네 가지로 함께 바꿨고, 하나만으로는 부족했다(표 \ref{t:e20}).

\begin{itemize}
\item \textbf{1차 내부 위상천이(EPS).} 1차 브리지 두 레그 사이에 위상차 D1·π를 주면 1차 전압 펄스 폭이 (1 − D1)·π가 된다\cite{b2}. 링크 5 V·11 A에서 첨두가 189 A에서 20 A로 줄어든다. 정상상태 파형의 정확한 계산으로 (전압비 r, 정규화 전류 j)마다 첨두가 가장 작은 (D1, 위상) 지령표를 만들었다(SPS(D1 = 0)는 해석식과 오차 없이 같음).
\item \textbf{전압비 정합 유지.} 버스를 V\textsubscript{W}에 붙잡지 않고 n·v\textsubscript{link} = v\textsubscript{bus}가 되도록 링크를 둔다(1008 V에서 링크 약 89 V, 버스 약 410 V). DAB 정격 밖 링크전력은 제동 저항이 받는다.
\item \textbf{첨두 상한.} 강등 상태에서 DAB 지령을 EPS 첨두 ≤ 1.5 × 정격점(R2-W460 68.7 A, R2-W410 90.0 A)이 되는 전류까지만 낸다.
\item \textbf{일시 감발과 이중 고장 전환.} H-브리지 고장이면 배터리 전압과 상관없이 400 V 부하 출력을 잠시 줄인 뒤 링크를 내린다. 소프트 투입이 5 ms 안에 끝나지 않으면 주 스위치를 강제로 닫지 않고 이중 고장(Stage A 협조 정지)으로 넘긴다(강제로 닫으면 693 A). 이중 고장의 최종 처리는 9.6절에서 정했다.
\end{itemize}

\Needspace{8\baselineskip}
{\footnotesize
\begin{longtable}{L{\dimexpr 0.3860\linewidth-2\tabcolsep\relax}L{\dimexpr 0.1316\linewidth-2\tabcolsep\relax}L{\dimexpr 0.1316\linewidth-2\tabcolsep\relax}L{\dimexpr 0.3509\linewidth-2\tabcolsep\relax}}
\caption{강등 방식별 DAB 첨두와 판정(E20, 평균모델 + 정확 해석)}\label{t:e20}\\
\toprule
\tabhead \textbf{방식} & \textbf{R2-W460} & \textbf{R2-W410} & \textbf{DAB 해석 첨두 최대(W460 / W410)} \\
\midrule
\endfirsthead
\toprule
\tabhead \textbf{방식} & \textbf{R2-W460} & \textbf{R2-W410} & \textbf{DAB 해석 첨두 최대(W460 / W410)} \\
\midrule
\endhead
\bottomrule
\endlastfoot
P0 기존 추천(버스 V\textsubscript{W} 유지, SPS) & 0/4 & 0/4 & 197.3 / 202.0 A \\
P1 P0 궤적 + EPS(상한 없음) & 0/2 & 0/2 & 135.7 / 136.4 A \\
P3n 정합 유지 + EPS 상한, 감발 없음 & 1/2 & 2/2 & 68.8 / 90.0 A. W460 930 V는 링크 61 V 정체 → 이중 고장 \\
\textbf{P3 P3n + 일시 감발(추천)} & \textbf{4/4} & \textbf{4/4} & \textbf{68.8 / 90.0 A} \\
\end{longtable}
}
\srcnote{주: 판정은 E17 고장 판정에 DAB 첨두(해석) ≤ 1.5 × 정격점을 더한 것이다. P3 감도 6종 중 W410 감쇠기 없음(경로 100.2 A)만 불통과.}

\begin{figure}[H]
\centering
\includegraphics[width=16.0cm]{\repo/validation/2026-10-24-dab-fault-modulation/E20/E20_dab_fault_modulation.png}
\caption{DAB 고장 모드 변조}\label{f:e20}
\srcnote{주: (a) 버스 504 V에서 링크전압에 따른 SPS·EPS 첨두, (b) 링크 5 V·11 A의 한 주기 인덕터 전류, (c) 1008 V·25 kW H-브리지 개방의 해석 첨두와 회로 첨두(TB09·TB10).}
\end{figure}

P3의 비용은 유지 중 버스가 V\textsubscript{W}보다 50–80 V 낮고(최저 373·337 V), 고장당 제동 에너지가 6.8 J에서 21.4 J로 늘며(분리 20 ms), V\textsubscript{A}가 상한 이하인 팩에서도 400 V 부하가 약 30 ms 감발된다는 것이다.

회로 확인(TB10)은 TB09 회로에 EPS 지령 블록과 게이트 생성기(1차 레그 2를 D1·T/2 지연)를 넣어 했고, 세 사례(1008·930 V H-브리지 개방, 930 V 켬)가 회로 판정을 통과했다. \textbf{DAB 링크 쪽 첨두는 192.0·178.1 A에서 79.4·75.3 A로 줄었다}(정격점의 1.73·1.64배). 링크는 0.49 V까지 붕괴해 TB09의 정체가 없어졌고, 소프트 투입 소자 에너지는 0.024·0.026 J(TB09 930 V는 1.22 J), 경로 순시 최대는 57.2·57.9 A였다. 다만 미리 정한 회로 기준(2배, 91.7 A)은 지켰으나 \textbf{해석 설계 목표 1.5배(68.7 A)는 링크 붕괴 구간에서 10–16 \% 넘었다.} 지령표 보간은 원인이 아니었다(상한에 붙인 보간 첨두 ≤ 1.0007배).



\section{회로 여유: 첨두 표, 전이 편류, 설계 상한(E21, TB11–TB13)}

TB10의 초과 원인을 확인하려고, 회로와 같은 소자 조건(누설 인덕터 1차 14.68 µH·20 mΩ, 이상 변압기 4.6:1, 데드타임 100 ns, 출력용량 1차 335 pF·2차 1,371 pF, 역방향 도통 2 V)으로 DAB 한 운전점의 주기 정상상태를 계산하는 스위칭 모의를 만들었다(E21). 소자 효과를 끈 모의는 정확 해석과 6개 운전점에서 0.10 \% 안으로 같았고, TB10 붕괴 구간(링크 ≥ 6 V)의 회로 첨두를 1008 V 150주기에서 최대 1.029·평균 1.003배, 930 V 136주기에서 최대 1.020·평균 1.001배로 재현했다(같은 주기의 회로/해석은 최대 1.131·1.098배). 효과별 기여는 버스 447 V·링크 52.8 V에서 해석 67.6 A → 데드타임 74.6 A → + 출력용량 74.0 A → + 역방향 도통 74.4 A, 버스 476 V·링크 23.6 V에서 68.0 → 64.6 → 72.3 → 72.3 A였다. 데드타임 100 ns가 붕괴 막바지의 1차 펄스 폭(125–150 ns)과 비슷한 크기라 펄스 폭과 위상이 실질적으로 바뀌기 때문이다.

모의가 회로 첨두를 평균 0.1–0.3 \%, 최대 2.9 \% 안에서 재현했으므로, 초과의 원인은 \textbf{데드타임과 그 동안의 출력용량 충방전}이다(역방향 도통의 기여는 작다). 대책으로 (D1, 위상) 지령표는 그대로 두고 \textbf{첨두 표만 같은 칸의 스위칭 정상상태 값}(버스 380·440·500 V 중 최댓값)으로 바꾸고, 설계 상한은 모의와 회로의 최대 차(2.9 \%)를 덮도록 1.5배 / 1.03 = 1.455배(R2-W460 66.7 A)로 두었다. 같은 표의 실제 링크전류로 평균모델도 계산했다(표 \ref{t:e21b}).

{\footnotesize
\begin{longtable}{L{\dimexpr 0.4035\linewidth-2\tabcolsep\relax}L{\dimexpr 0.2982\linewidth-2\tabcolsep\relax}L{\dimexpr 0.2982\linewidth-2\tabcolsep\relax}}
\caption{데드타임 반영 첨두 표의 확인(E21, 평균모델 P3s)}\label{t:e21b}\\
\toprule
\tabhead \textbf{시험} & \textbf{R2-W460} & \textbf{R2-W410} \\
\midrule
\endfirsthead
\toprule
\tabhead \textbf{시험} & \textbf{R2-W460} & \textbf{R2-W410} \\
\midrule
\endhead
\bottomrule
\endlastfoot
첨두 표로 정한 지령의 실제 첨두(버스 337–505 V × 링크 1–95 V, 각 273점) & 최대 1.466배(≤ 1.5배) & 최대 1.462배(≤ 1.5배) \\
기본 사례(930·1008 V / 860·920 V × 5·25 kW) & 4/4, DAB 첨두 66.8 A, 경로 ≤ 71.2 A & 4/4, 87.4 A, 경로 ≤ 66.4 A \\
버스 최저 / 제동 에너지(분리 20 ms) / 소프트 투입 & 373 V / ≤ 19.2 J / ≤ 0.029 J & 342 V / ≤ 5.6 J / ≤ 0.025 J \\
링크 붕괴 시간 & ≤ 10.0 ms(시간 상한에 닿음) & ≤ 7.6 ms \\
감도(감쇠기 없음, V\textsubscript{D} 1.4·3.0 V, 분리 50·3 ms) & 5/5 & 4/5(감쇠기 없음 100.3 A, 기존과 같음) \\
\end{longtable}
}
\srcnote{주: 판정은 E17 고장 판정 + 스위칭 반영 DAB 첨두 ≤ 1.5 × 정격점(68.7·90.0 A).}

회로 확인은 같은 회로·지령 블록으로 세 단계에 걸쳐 새 테스트벤치로 했다. 판정 기준은 시험 전에 TB10의 2배에서 \textbf{1.5배(68.7 A)}로 좁혀 고정했다(표 \ref{t:tb1013}). TB11은 첨두 표 값과 상한만, TB12는 게이트 생성기만, TB13은 상한 값만 바꿨다.

\Needspace{12\baselineskip}

{\footnotesize
\begin{longtable}{L{\dimexpr 0.28\linewidth-2\tabcolsep\relax}L{\dimexpr 0.18\linewidth-2\tabcolsep\relax}L{\dimexpr 0.18\linewidth-2\tabcolsep\relax}L{\dimexpr 0.18\linewidth-2\tabcolsep\relax}L{\dimexpr 0.18\linewidth-2\tabcolsep\relax}}
\caption{TB10–TB13 회로 모델의 판정(R2-W460, 분리·감발 3 ms. 각 칸은 사례 A 1008 V / B 930 V, H-브리지 개방)}\label{t:tb1013}\\
\toprule
\tabhead \textbf{항목} & \textbf{TB10} & \textbf{TB11} & \textbf{TB12} & \textbf{TB13} \\
\midrule
\endfirsthead
\toprule
\tabhead \textbf{항목} & \textbf{TB10} & \textbf{TB11} & \textbf{TB12} & \textbf{TB13} \\
\midrule
\endhead
\bottomrule
\endlastfoot
바꾼 것 & EPS 지령 블록·게이트 도입 & 첨두 표(스위칭 값)·상한 & 게이트 생성기\newline(편류 방지) & 설계 상한 \\
설계 상한 & 1.5배\newline(68.7~A) & 1.455배\newline(66.7~A) & 1.455배\newline(66.7~A) & \textbf{1.387배\newline(63.6~A)} \\
DAB 첨두 판정 기준 & \textbf{2배\newline(91.7~A)} & \textbf{1.5배\newline(68.7~A)} & 1.5배\newline(68.7~A) & 1.5배\newline(68.7~A) \\
회로 판정 & 통과 / 통과 & \textbf{불통과 / 불통과} & \textbf{불통과} / 통과 & 통과 / 통과 \\
1.5배 기준 통과 사례 & 해당 없음 & \textbf{0/2} & \textbf{1/2} & \textbf{2/2} \\
DAB 첨두: 유지(A) & 58.5 / 62.7 & 58.5 / 62.7 & 58.0 / 58.0 & 58.0 / 58.0 \\
DAB 첨두: 붕괴(A) & 79.4 / 75.3 & 88.5 / 68.9 & 70.0 / 68.5 & 66.4 / 64.4 \\
DAB 첨두: 소프트 투입(A) & 17.6 / 17.5 & 22.1 / 17.6 & 22.2 / 17.6 & 58.3 / 19.2 \\
1.5배 초과 주기 & 1082 / 916 & 7 / 3 & 3 / 0 & 0 / 0 \\
강등 상태 주기 & – & 2801 / 2688 & 2801 / 2694 & 2801 / 2801 \\
1차 첨두(A) & 17.3 / 16.4 & 19.2 / 15.0 & 15.2 / 14.9 & 14.4 / 14.0 \\
경로전류 단독 최대(A) & 57.2 / 57.9 & 57.2 / 57.9 & 57.0 / 57.9 & 57.0 / 57.9 \\
경로·필터 중 최댓값(A) & 57.3 / 57.9 & 57.3 / 57.9 & 57.1 / 57.9 & 57.1 / 57.9 \\
버스 최저(V) & 409.6 / 380.6 & 409.6 / 380.6 & 409.5 / 380.4 & 409.5 / 380.4 \\
소프트 투입 진입 링크(V) & 0.49 / 0.49 & 1.06 / 0.49 & 1.07 / 0.49 & 4.27 / 0.68 \\
소프트 투입 소자 에너지(J) & 0.024 / 0.026 & 0.028 / 0.026 & 0.028 / 0.026 & 0.048 / 0.027 \\
제동 저항 에너지(J) & 6.9 / 5.8 & 7.6 / 6.2 & 7.6 / 6.3 & 8.1 / 6.7 \\
강등 완료(ms) & 13.80 / 13.18 & 14.51 / 13.95 & 14.51 / 13.98 & 14.51 / 14.51 \\
\end{longtable}
}
\srcnote{주: TB10은 2배 기준 판정이며, TB11–TB13의 판정은 TB10의 회로 한계에 강등 상태(9–11)의 DAB 링크 쪽 순시 첨두 ≤ 1.5 × 정격점(68.7 A)을 더한 것이다. DAB 첨두는 링크 쪽 순시 첨두다. 경로전류 단독은 순시 |i\textsubscript{p}|의 최대, 경로·필터 중 최댓값은 순시 |i\textsubscript{p}|와 |i\textsubscript{h}| 가운데 큰 값이다. 버스 최저는 5 µs 평균, 제동 저항 에너지는 분리 3 ms, 강등 완료는 사건 0.5 ms 기준 시각이다.}

\textbf{TB11은 두 사례 모두 1.5배 기준에서 불통과였다.} 다만 실패의 성격이 TB10과 다르다. 1.5배를 넘은 주기가 TB10의 1,082·916주기에서 7·3주기로 줄었고, 정상 주기의 회로 첨두는 같은 (버스, 링크)의 스위칭 모의와 0.3 A 안에서 같았으며, 넘은 주기는 모두 지령표의 칸이 바뀌는 순간이었다. 사례 A의 5.215 ms에서는 EPS 지령이 (위상, D1) = (−0.029 rad, 0.278)에서 (0.199 rad, 0.379)로 한 주기 만에 옮겨 갔고 그 주기의 첨두가 88.5 A였다(앞 주기 66.9 A, 두 지령의 정상상태 첨두는 66.6·66.5 A). 즉 \textbf{남은 초과는 정상상태가 아니라 지령 전이의 과도}다. 평균모델(E21)은 이 과도를 다루지 않으므로, 상태 전이 시각은 회로와 0.02 ms 안에서 맞았지만 첨두 판정은 회로에서만 드러났다.

게이트 생성기만 떼어 내 같은 지령 계단을 이상 스위칭으로 계산하면(버스 432 V·링크 67.5 V) 계단 전 70.1 A이던 첨두가 계단 뒤 1–4주기에 92.6·76.9·76.8·76.7 A가 되고 누설 전류에 −16.6 A의 직류 편류가 지속되었다(이상 모델에는 데드타임 감쇠가 없어 편류가 R/L(0.73 ms)로만 줄어든다. 회로의 TB11 A에서는 계단 뒤 첨두가 88.5 → 69.4 → 68.2 A로 수 주기 안에 줄었다). SPS 위상 변화에는 전이 반주기에 이전·새 위상의 평균을 쓰는 직류 편류 방지가 있었지만 EPS의 레그 2 지연(D1)은 상향 에지에서 한 주기 동안 그대로 잡았기 때문이다. 구형파 레그의 에지가 Δ만큼 옮겨질 때 첫 에지를 Δ/2만 옮기면 직류 편류가 남지 않으므로 이를 레그 2에도 적용하고, 큰 칸 전환은 위상 반주기당 0.01 rad, D1 에지당 0.005로 나눠 옮기도록 고쳤다(TB12). 고친 생성기는 같은 계단에서 70.6·69.6·69.7·68.6 A, 편류 ≤ 2.7 A였다. 이 변화율 상한은 정상 붕괴 중의 지령 변화(주기당 D1 약 0.0006)보다 충분히 커서 평소 동작에는 영향이 없다.

\textbf{TB12는 사례 B만 통과했다(1/2).} 사례 A는 5.26 ms에 시작한 한 전이의 끝 세 주기에서 최대 70.0 A로 기준을 1.3 A 넘었다. 전이 경로 중간의 (D1, 위상)은 정상 상한 칸보다 첨두가 조금 높고 D1 변화의 직류 편류가 몇 A 남기 때문이며, 전이 주기 첨두는 정상 주기 상한(66.7 A) 위로 최대 4.9 \%였다. 상한 탐색을 버스·링크 전압비에 연속이 되게 하는 안은 TB11 사례 A의 궤적에서 오프라인으로 점검했지만 주기당 지령 도약이 11 \%(위상 0.227 → 0.202 rad)만 줄었다(같은 전압비에서 이웃 칸의 최적 (D1, 위상)이 최대 0.41 rad 차이). 그래서 TB13은 전이 첨두를 여유로 넣어(4.9 \% → 5 \%) 설계 상한을 1.5배 / (1.03 × 1.05) = \textbf{1.387배(63.6 A)}로 낮췄다. 평균모델(E21 P3t)에서는 기본 사례가 모두 통과했고(R2-W460 4/4, R2-W410 4/4, 스위칭 반영 첨두 63.7·83.4 A), 링크가 덜 붕괴해 소프트 투입 소자 에너지가 R2-W460에서 최대 0.058 J(P3s 0.029 J)로 늘었지만 TB09 이전 방식(1.22 J)보다 한 자릿수 작다.

\textbf{TB13은 두 사례 모두 1.5배 기준을 포함한 회로 판정을 통과했다.} DAB 링크 쪽 첨두는 66.4·64.4 A로 정상 주기 상한(63.6 A) 위로 최대 4.4 \%, 기준까지 여유는 2.4 A다. 강등 중 DAB 첨두는 해석(최적 지령표) → 스위칭 정상상태(데드타임·출력용량, 3 \%) → 지령 전이(변화율 상한과 여유 5 \%)의 세 단계를 거쳐 회로에서 설계 목표 안에 들었다(그림 \ref{f:e21}). 대가는 붕괴 구간의 DAB 전류가 약 5 \% 줄어 링크가 덜 내려간 것이다. 사례 A는 진입 링크 4.27 V로 소자 에너지가 0.048 J로 늘고 소프트 투입 중 첨두도 58.3 A가 되었으며, 사례 B는 붕괴가 시간 상한에 닿아 강등 완료가 14.51 ms로 늦어졌다.


\begin{figure}[!htb]
\centering
\includegraphics[width=16.0cm]{\repo/validation/2026-10-25-dab-switching-margin/E21/E21_dab_switching_margin.png}
\caption{데드타임 반영 첨두 표(E21)와 회로 확인(TB10–TB13)}\label{f:e21}
\srcnote{주: (a) 첨두 상한이 허용하는 DAB 전류(이상 표 1.5배, 스위칭 표 1.455배·1.387배), (b) TB10 사례 A의 붕괴 구간 첨두: 해석·스위칭 모의·회로, (c) 1008 V·25 kW H-브리지 개방의 링크전압과 회로 DAB 첨두(TB10–TB13).}
\end{figure}

\section{보호 상태기계와 SALS 연동(E29, TB23)}

9.1–9.4절의 강등 순서는 “400 V 부하가 명령 20 ms 뒤 5 kW로 줄어든다”는 가정 위에 있었다. 실제로는 V\textsubscript{A}가 부품 상한을 넘는 팩은 접촉기로 부하를 분리하고, 상한 안의 팩은 SALS 메시지로 일시 감발을 전한다. 메시지는 늦거나 손실될 수 있고, 부하는 계약 변화율로만 줄며, M0 등급(Stage C·V2L)은 줄지 않는다. 메시지 지연·손실·재송신·워치독은 해석식으로 계산했다. 이 응답을 평균모델 상태기계(강등 순서 E21 P3t)에 넣어 다시 확인했고, 연동 설계는 다음과 같다.

\begin{itemize}
\item 고장을 검출하면 SALS 감발(D 등급 허용치 0, M1 한도 0)과 Stage C 회생 차단을 함께 낸다. V\textsubscript{A}가 상한을 넘는 팩은 접촉기 분리도 바로 명령한다.
\item 감발만 하는 팩에서는 H-브리지 고장 유지(상태 9)를 시간이 아니라 측정으로 끝낸다. 기준은 부하전류 추정 i\textsubscript{p} − (C\textsubscript{b} + C\textsubscript{dmp})·dv\textsubscript{bus}/dt이고, 문턱 I\textsubscript{COL}은 22.2·28.9 A(R2-W460·R2-W410)다. 30 ms 안에 부하가 줄지 않으면 접촉기 분리로 넘어간다.
\item 이중 고장(상태 13)이면 Stage A를 멈추고, 400 V 부하 분리·감발과 링크 제동 저항의 과전압 보호를 함께 쓴다.
\end{itemize}

\Needspace{8\baselineskip}
{\footnotesize
\begin{longtable}{L{\dimexpr 0.28\linewidth-2\tabcolsep\relax}L{\dimexpr 0.30\linewidth-2\tabcolsep\relax}L{\dimexpr 0.42\linewidth-2\tabcolsep\relax}}
\caption{보호 상태기계–SALS 연동의 결과(E29, 평균모델)}\label{t:e29}\\
\toprule
\tabhead \textbf{시험} & \textbf{조건} & \textbf{결과} \\
\midrule
\endfirsthead
\toprule
\tabhead \textbf{시험} & \textbf{조건} & \textbf{결과} \\
\midrule
\endhead
\bottomrule
\endlastfoot
연동 강등 & 팩 4개 × 부하 구성 3개 × 메시지 시나리오 6개 + 접촉기 50 ms 감도 2 & \textbf{74/74}. 감발만 하는 팩(930·832 V, 25 kW, 지연 5 ms)의 강등 완료 23–30 ms, 제동 5.8–10.0 J(20 ms 계단 가정 31.5 ms, 22.1 J) \\
메시지 손실 & 감발 메시지 두 번 이상 손실, 컴프레서 메시지 전부 손실 & 접촉기 분리로 넘어감. 유지 50 ms, 제동 최대 53.3 J(기준 64 J) \\
48 V 회생 & 고장 순간 순부하 −3.2 kW(Stage C 회생) & 회생 차단 없음 0/4(버스 432.7–527.0 V), 차단 있음 12/12. 차단 시작은 검출 뒤 R2-W460 ≤ 1.0 ms, R2-W410 ≤ 0.5 ms(4 kW/ms) \\
감발 뒤 남는 부하 & M0 9.1 kW 남음 & 소프트 투입 소자 에너지 0.75 J(5 kW 가정의 24배). 13 kW까지 통과(2.7 J), 15.1 kW는 시간 상한에 걸려 이중 고장으로 \\
이중 고장 & Stage A 정지 + 부하 분리 + 제동 유지(C\textsubscript{A} 400 µF) & 링크 ≤ 110.1 V, 직렬 ≤ 114.2 V. 경로전류가 1 A 아래까지 4.7–7.3 ms로 5 ms 기준 1/4(9.6절에서 설계 변경) \\
\end{longtable}
}
\srcnote{주: 이 표는 150 kHz Stage A(C\textsubscript{A} 400 µF)로 한 값이다. A25R에서 다시 계산한 이중 고장과 설계 변경은 9.6절에 있다. 자료: 검증 근거 E29.}

\textbf{새로 찾은 위험은 48 V 회생이다.} 고장 순간 버스 순부하가 음이면 다이오드 정류 유지의 극성이 바뀌어 버스가 올라, Stage C 회생 차단을 설계 요구로 두었다(10장). \textbf{남는 부하의 대가는 소프트 투입 에너지다.} M0 부하가 남아 링크가 덜 붕괴하면 바이패스 소프트 투입 소자의 에너지가 커지며, 이는 9.6절에서 풀었다.

회로 확인(TB23)은 TB13 회로에 같은 부하 블록과 상태기계 확장을 넣어 했다(SALS 시간 0.15배, 기준값도 같은 시간). 다섯 사례의 판정이 모델과 모두 같았고(통과 4, 회생 차단을 뺀 사례 불통과, 버스 481.0 V), 상태 전이 시각은 0.065 ms 안, 소프트 투입 에너지는 0.476 J(모델 0.471 J)였다.

\textbf{H4와의 관계.} 감발 메시지가 늦거나 손실되어도 강등 순서가 성립하므로(74/74, 표 \ref{t:e29}), 연동은 축약 모델에서 확인되었다. 전압 의존 예산과 전기 모델 결합의 일반화는 남은 과제다(18장).

\section{바이패스 소프트 투입과 이중 고장의 설계 변경(E37, E38, TB31)}

A25R로 바꾸고 부품 정격을 데이터시트로 판정하자(16장) 두 가지가 불통과로 남았다.

\begin{itemize}
\item \textbf{바이패스 소프트 투입 소자(P6).} SALS 연동에서 남는 부하 때문에 링크 붕괴가 시간 상한에 걸리면 직렬전압이 남은 채(최대 −25.7 V) 소프트 투입을 시작한다. 최악 에너지 0.749 J에서 바이패스 GaN(EPC2305) 1개의 접합온도는 176 °C였고(E34, 데이터시트 과도 열임피던스), 선형 영역용 MOSFET 5종은 펄스 열임피던스 값이 없어 보수 계산으로 245.5 °C 이상이었다(E37).
\item \textbf{이중 고장 경로전류(F2).} C\textsubscript{A}가 2.4 mF로 커지자 Stage A 정지 뒤 V\textsubscript{A} 노드의 에너지가 6배가 되어, 직렬전압이 최대 123.3 V(3사례 120 V 초과), 경로전류 5 ms 기준 0/4(7.6 ms 이상)가 되었다(E36). V\textsubscript{A} 노드 방전 저항 1 Ω은 2/4였고, 방전 첨두 430–504 A로 4사례 중 3사례가 클램프 스위치 한도(450 A)를 넘었다(E37). 원인은 연결된 400 V 부하가 V\textsubscript{A} 노드와 버스 커패시터에서 전력을 빼며 버스를 저전압 차단(250 V)까지 끌어내리는 동안 경로전류가 남는 것이다.
\end{itemize}

두 문제를 부하 쪽 설계 변경으로 풀었다(E38, 판정 기준은 E29와 부품 정격 표 그대로).

\textbf{Stage C 보호 일시 정지.} H-브리지 고장 강등의 링크 붕괴와 소프트 투입 동안 Stage C를 4 kW/ms로 0까지 줄여 멈추고, 강등 0.5 ms 뒤 다시 켠다. 허브 내부의 제어이고, 멈춤 시간(정지 램프 시작부터 재개 램프 끝까지)은 8.7–12.4 ms다. 에너지는 포트별로 다르다(E46 정정): 고압 버스 쪽 Stage C 입력은 사각 상계 5.3 kW × 12.41 ms = 65.8 J(E38, 같은 4 kW/ms 램프의 공급 결손 적분은 58.8 J)이고, 48 V 쪽은 별도 시험(E43)에서 +5 kW 부하를 12.4 ms 완전히 멈춘다고 두어 부하 62 J, 내부저항 20 mΩ 세 점의 배터리 방전 64.7–65.3 J였다(75.5 J는 OCV 44 V·60 mΩ에서 35.6 V로 기준을 못 지킨 행). 이 동안 48 V는 Stage C가 조정하지 않고 배터리 전압을 따르며, 무제한 운전 범위(36–52 V)를 지키는 배터리 요구는 10장에 도출했다(E43). 결과는 연동 강등 74/74(지정한 H-브리지 개방·메시지 시나리오. 반전압이 부품 상한 밖인 1008·920 V 38건은 바로 분리, 상한 안인 930·832 V 36건은 감발 경로이며 그중 25건은 접촉기 분리 없이 400 V 부하를 복원하고 30 ms 안에 감발을 확인하지 못한 11건만 분리로 바뀜, E46), 시간 상한 투입 30사례 → 0사례, 최악 소프트 투입 에너지 0.036 J·첨두 104 W, EPC2305 1개의 접합온도 108.0 °C(통과, 16장의 보수 규칙으로 109.6 °C), 바이패스 전류 최대 45.0 A(소자 1개 선형 한도 164.5 A, E46)다. 강등 완료는 0.5–2.3 ms 빨라졌고 제동 에너지 최대는 51.9 J다.

\textbf{이중 고장의 400 V 부하 정지.} 이중 고장에 들어가면 Stage A를 멈추고, 전용 배선 신호로 모든 400 V 부하(Stage C 포함)를 정해진 램프로 0까지 줄인다. V\textsubscript{A} 노드 방전 저항은 두지 않는다. 램프 시간을 바꿔 가며 계산한 결과는 표 \ref{t:f2}과 같다.

{\footnotesize
\begin{longtable}{L{\dimexpr 0.16\linewidth-2\tabcolsep\relax}L{\dimexpr 0.12\linewidth-2\tabcolsep\relax}L{\dimexpr 0.22\linewidth-2\tabcolsep\relax}L{\dimexpr 0.18\linewidth-2\tabcolsep\relax}L{\dimexpr 0.16\linewidth-2\tabcolsep\relax}L{\dimexpr 0.16\linewidth-2\tabcolsep\relax}}
\caption{이중 고장의 400 V 부하 정지 램프별 결과(E38, 평균모델, A25R)}\label{t:f2}\\
\toprule
\tabhead \textbf{정지 램프} & \textbf{판정} & \textbf{경로전류 ≤ 1 A까지} & \textbf{직렬 최대} & \textbf{버스 끝} & \textbf{H-브리지 역방향 도통 소자} \\
\midrule
\endfirsthead
\toprule
\tabhead \textbf{정지 램프} & \textbf{판정} & \textbf{경로전류 ≤ 1 A까지} & \textbf{직렬 최대} & \textbf{버스 끝} & \textbf{H-브리지 역방향 도통 소자} \\
\midrule
\endhead
\bottomrule
\endlastfoot
1 ms & 4/4 & 0.62–0.63 ms & 64.6 V & ≥ 395 V & 100 °C \\
2 ms & 4/4 & 1.18–1.27 ms & 79.4 V & ≥ 379.6 V & 100 °C \\
\textbf{3 ms(선정)} & \textbf{4/4} & \textbf{1.83–1.98 ms} & \textbf{94.8 V} & ≥ 363 V & 100 °C \\
5 ms & 4/4 & 3.20–4.89 ms & 114.0 V & ≥ 328.9 V & 103 °C \\
\end{longtable}
}
\srcnote{주: 5 ms도 4/4지만 직렬전압이 비교기 문턱(100 V)을 넘고, R2-W410 920 V는 다이오드 도통 직전(114 V)이다. 그래서 요구값을 3 ms로 정했다. 자료: 검증 근거 E38.}

이중 고장 뒤 V\textsubscript{A} 노드는 V\textsubscript{bat}/2 근처(428–503 V)로 남아 재기동 사전충전 부담이 가볍다. 서비스는 바뀌지 않는다(이전에도 버스가 4–6 ms 안에 저전압 차단에 닿아 400 V 부하가 끊겼다).

회로 확인(TB31)은 TB23 회로(R2-W460)에 C\textsubscript{A} 2.4 mF, Stage C 일시 정지, 400 V 부하 정지를 넣어 했고 두 사례 모두 통과했다. 930 V·25 kW H-브리지 개방(메시지 지연 5 ms)의 소프트 투입 소자 에너지는 0.0227 J(모델 0.0214 J), EPC2305 1개 105.0 °C(모델 104.7 °C)였고, 1008 V·25 kW 이중 고장에서는 경로전류가 Stage A 정지 1.110 ms 뒤 1 A 아래로 떨어졌다(모델 1.113 ms, 기준 5 ms). 상태 전이 시각은 모델과 0.035 ms 안이었다.

이 변경으로 보호 순서의 불통과는 없어졌지만 \textbf{새 인터페이스 요구 두 가지}가 생겼다. Stage C 일시 정지(허브 내부, 최대 12.4 ms)는 허브 안에서 구현할 과제이고(Stage C 제어기의 입력 제한은 아직 구현하지 않았다, 19장), 이중 고장 때 400 V 부하 정지(전용 배선 신호 뒤 ≤ 3 ms)는 컴프레서·히터 등의 실제 응답에 달려 있어 차량 수준에서 확인해야 한다(19장).
